\documentclass[11pt]{article}
\usepackage[margin=1in]{geometry}
\usepackage{amsmath,amssymb,amsthm,mathtools}
\usepackage{hyperref}
\newtheorem{theorem}{Theorem}[section]
\newtheorem{lemma}[theorem]{Lemma}
\newtheorem{proposition}[theorem]{Proposition}
\newtheorem{remark}[theorem]{Remark}
\newcommand{\OO}{\mathcal O}
\newcommand{\CC}{\mathcal C}
\newcommand{\DD}{\mathcal D}
\newcommand{\Jac}{\operatorname{Jac}}
\newcommand{\Bl}{\operatorname{Bl}}
\newcommand{\Fit}{\operatorname{Fit}}
\title{Isolated positive-characteristic singularities:\\
canonical-module flattening and uniform obstacle envelopes}
\author{Research continuation 377: working proof for independent verification}
\date{8 October 2026}
\begin{document}
\maketitle
\begin{abstract}
We propose an envelope-continuity argument for a normal isolated algebraic
singularity over a perfect field of positive characteristic, in arbitrary
dimension. The construction uses a blowup of the first Fitting ideal of a
canonical module, rather than a smooth resolution. On the resulting normal
primary model an effective exceptional boundary makes the log canonical
divisor Cartier. Relative Serre vanishing gives a uniform generation statement
for asymptotic test ideals of this pair. Jacobian-corrected subadditivity then
gives an error of order $1/m$ on every normalized centered semivaluation,
including those with nonzero support. This is a working proof, not an accepted
correctness grant. It addresses a positive-characteristic isolated normal
sector of P12.3; the full excellent and nonnormal original problem stays open.
\end{abstract}

\section{Objects and the precise proposed conclusion}
Let $k$ be perfect of characteristic $p>0$, let $B$ be a normal integral
affine $k$-variety, and let $x$ be a closed point of codimension
$d=\dim\OO_{B,x}\geq2$. After shrinking $B$, suppose $B\setminus\{x\}$ is
regular. Fix an $x$-primary ideal $\mathfrak a$. Let
$S_{\mathfrak a}$ be the compact section of centered semivaluation rays
normalized by $v(\mathfrak a)=1$. We use additive valuations; values of
individual functions may be $+\infty$. A primary ideal has finite value
at every point of this section.

The geometric comparison available from the qualified P12.2 paper is used
in its exact scope: primary-ideal profiles
$-v(I)/N$ have independent negative effective Cartier, exceptional-nef
models; their uniform closure is the continuous geometric cone $\CC$;
the geometric decreasing class $\DD$ consists of the appropriate limits
of one decreasing net from $\CC$. No envelope-continuity, upward-limit,
energy, or resolution assertion is imported. In particular,
$\CC$ contains zero and the negative constants, is closed under maxima,
and permits a negative constant shift. We retain its nonpositive convention.
For a finite continuous obstacle $\phi$ define the raw envelope
\[
 P\phi(v)=\sup\{u(v):u\in\DD,\ u\leq\phi\}.
\]
Constants make this family nonempty and bounded below.

\begin{theorem}[Proposed isolated normal sector]\label{thm:main}
Under these hypotheses $P\phi$ is continuous and belongs to $\CC$ for every
finite continuous $\phi$ on $S_{\mathfrak a}$. It is the greatest geometric
minorant and
$\|P\phi-P\psi\|_\infty\leq\|\phi-\psi\|_\infty$.
No smooth-model hypothesis is imposed. The same construction gives a
greatest $\DD$ minorant for every bounded upper semicontinuous obstacle,
as the limit of one decreasing net of continuous geometric envelopes.
\end{theorem}

The working proof below is written so that canonical flattening, actual
primary-model realization, trace generation, and all-support estimates can
each be independently checked. It does not claim this theorem for a
nonisolated singularity, mixed characteristic, an imperfect field, an arbitrary
nonalgebraic excellent local ring, or a nonnormal germ.

\section{Flattening and a Cartier log canonical divisor}
\begin{lemma}\label{lem:flatten}
Every finite collection of primary-ideal profiles is represented on a normal
actual primary blowup $Y\to B$, isomorphic away from $x$, for which there is
an effective integral Weil divisor $\Delta$, supported on the fiber, with
$L=K_Y+\Delta$ Cartier. There is an effective full-fiber Cartier divisor $F$
with $A=\OO_Y(-F)$ ample and globally generated, and $Y=\Bl_J B$ for an
actual $x$-primary ideal $J$ satisfying $J\OO_Y=A$.
\end{lemma}
\begin{proof}
Include the gauge ideal in the finite collection, then blow up their product
and normalize.
All the ideals become invertible: over a local domain, if the product of
nonzero ideals is invertible, each factor is invertible, by the fractional
inverse obtained from the other factors. Normalization is finite because
the schemes are excellent. The pulled-back product ideal defines a full-fiber
effective Cartier divisor $F_0$, with $\OO(-F_0)$ ample over the affine base.

On this normal scheme the canonical module $\omega$ is coherent, torsion
free of rank one, and invertible away from the fiber. Set $G=\Fit_1(\omega)$.
This ideal is the unit ideal away from the fiber. The strict transform of
$\omega$ on $\Bl_G$ is invertible, by Stacks 15.27.3--4. Explicitly, on an
affine blowup chart $R'=R[G/a]\subset R_a$, its quotient by $a$-power
torsion has $\Fit_1=R'$; its $\Fit_0$ vanishes after inverting $a$ and hence
vanishes in $R'$. The Fitting criterion makes it locally free of rank one.
This argument includes charts where some individual presentation generators
vanish; it does not assume that $\omega$ was initially invertible.

Normalize this blowup. The pullback of that invertible strict transform
is still invertible. Call it $\mathscr L$. On the puncture it agrees with
$\omega_Y$. A rational trivialization gives a Cartier divisor $L_0$
representing $\mathscr L$, and $K_Y-L_0$ is an integral Weil divisor supported
on the fiber. The pullback of $F_0$ has strictly positive coefficient at
each prime divisor lying in the fiber. A sufficiently large multiple
$rF_0$ therefore dominates $K_Y-L_0$. Put
\[
 \Delta=rF_0-(K_Y-L_0)\geq0,
 \qquad K_Y+\Delta=L_0+rF_0.
\]
Thus the log canonical divisor is Cartier without making $Y$ regular.

Here and below the effective exceptional divisor from the Fitting blowup
may be zero if $G=\OO$. Let it be $E_G$. For sufficiently large $b$,
$\OO(-bF_0-E_G)$ is ample over $B$: it is the relative ample blowup line
bundle tensor a sufficiently high pullback of an ample line bundle on the
previous scheme. Finite normalization preserves ampleness. Its divisor has
full-fiber support. Replace this line bundle by a sufficiently high power
which is generated by global sections and very ample over $B$.

If its effective divisor is $F$, set $J=\Gamma(Y,\OO(-F))$. Normality of
$B$ and codimension at least two give
$\Gamma(Y\setminus Y_x,\OO)=\Gamma(B,\OO)$, so $J$ is an ideal of $B$.
It is unit off $x$ and contains a power of $\mathfrak a$, by full-fiber
support and coherence. It is proper: its generated pullback vanishes at
every point of the nonempty fiber. Hence $J$ is primary.

The evaluation map gives $J\OO_Y=\OO(-F)$ and a closed immersion into
$\mathbf P_B(J)$ by very ampleness. On the puncture its image is the graph
defined by $J$. Since $Y$ is integral, that image is the scheme-theoretic
closure of this graph. The closure is exactly $\operatorname{Proj}
\bigoplus_{m\geq0}J^m=\Bl_J B$. This proves actual primary realization,
not just domination by a projective modification. Replacing $F$ by the
chosen power changes neither its support nor the Cartier property of
$K_Y+\Delta$.
\end{proof}

\begin{remark}
Two bounded corrections will be used. Since $Y$ is smooth over the perfect
field off the fiber, $\Jac(Y/k)$ is unit there, so
$J^a\OO_Y\subset\Jac(Y/k)$ for some integer $a$.
Also $\Delta\leq bF$ for some integer $b$: both are exceptional and
$F$ has positive coefficient on every exceptional prime. Consequently, for
every $v\in S_{\mathfrak a}$,
\[
 0\leq v(\Jac(Y/k))\leq a\,v(J),\qquad
 v(bF)=b\,v(J),\qquad v(J)\leq N_J
\]
for some $N_J$, since $\mathfrak a^{N_J}\subset J$.
The first inequality concerns an actual coherent ideal at the center of $v$;
it is valid even if the semivaluation has nonzero kernel.
\end{remark}

\section{Sections and the raw model-obstacle envelope}
Let $C$ be an integral Cartier divisor on $Y$ supported on the fiber. Its
function $f_C$ at a semivaluation is the value of a local equation for $C$.
This is independent of the equation since units have value zero. Properness
gives a center for every normalized semivaluation; a power of the primary
gauge contained in $J$ ensures a nonzero blowup denominator. The resulting
model function is finite and continuous on the whole section.

Define
\[
 I_m=\Gamma(Y,\OO_Y(mC)),\qquad
 \mathfrak b_m=I_m\OO_Y\cdot\OO_Y(-mC).
\]
Hartogs extension on $B$ identifies $I_m$ with an ideal in $B$.
It is unit off $x$, contains a primary power, and is either primary or $B$.
The nonzero ideals $\mathfrak b_m$ are graded, and
$\mathfrak b_m\OO(mC)$ is generated by the finitely many generators of
$I_m$. At every centered semivaluation,
\[
 u_m=-\frac{v(I_m)}m
     =f_C-\frac{v(\mathfrak b_m)}m\leq f_C.
\]
The value of every $I_m$ is finite. The subadditivity of $v(I_m)$ gives
the all-index Fekete limit
\[
 u=\lim_{m\to\infty}u_m=\sup_m u_m.
\]
We do not assert monotonicity of the sequence $u_m$.

\begin{lemma}\label{lem:raw}
The raw $\CC$ and $\DD$ envelopes of the continuous obstacle $f_C$
both equal $u$. This equality alone does not assert positivity of $u$.
\end{lemma}
\begin{proof}
Each $u_m$ is a primary-ideal profile or zero, hence an admissible continuous
minorant. If $h=-v(I)/q$ is an elementary profile below $f_C$, comparison at
the prime divisors of the normal $Y$ gives
$I^r\subset I_{rq}$ for every $r$: the inequality says that each section
has nonnegative order as a section of $\OO(rqC)$, and the height-one
intersection criterion on a normal scheme gives regularity everywhere.
Exceptional prime divisors give normalized centered divisorial tests;
at prime divisors off the fiber the orders of the sections are already
nonnegative because the sections belong to $B$ and $C$ is zero there.
Therefore $h\leq u_{rq}\leq u$.

For a general $h\in\CC$ below $f_C$, uniformly approximate it by elementary
profiles. A negative rational constant shift converts an approximation
within $\varepsilon$ into an elementary profile below $f_C$; this shift is
realized by multiplying a sufficiently divisible power by a power of
$\mathfrak a$. Let $\varepsilon\to0$.

Finally, if $h\in\DD$ is below $f_C$, write $h_\alpha\downarrow h$ with
$h_\alpha\in\CC$. For every $\varepsilon>0$, compactness applied to the
decreasing closed sets
$\{h_\alpha\geq f_C+\varepsilon\}$ gives one index with
$h_\alpha<f_C+\varepsilon$ everywhere. Then
$h_\alpha-\varepsilon\in\CC$ is below $f_C$, and $h\leq u+\varepsilon$.
This also covers $h$ with some value $-\infty$. It proves equality of both
raw suprema before the continuity argument.
\end{proof}

\section{A relative generation argument on a singular model}
All test ideals in this paper are big, nonfinitistic test ideals. Their
formation commutes with localization. For a graded family $\mathfrak b_\bullet$
and $m\geq1$ write
\[
 t_m=\tau(Y,\Delta,\mathfrak b_\bullet^m),\qquad
 s_m=\tau(Y,\mathfrak b_\bullet^m).
\]
An asymptotic test ideal is the unique maximal member among
$\tau(Y,\Delta,\mathfrak b_\ell^{m/\ell})$.
The family is directed by common multiples, using power invariance of test
ideals and gradedness, and its union stabilizes by noetherianity.
Simultaneous stabilization of finitely many such families uses one common
multiple. No smooth Frobenius-root definition is used here.
Power invariance itself follows from the big trace formula: the exponents
$r\lceil qt/r\rceil$ and $\lceil qt\rceil$ differ by at most $r-1$.
One inclusion is immediate; for the reverse inclusion multiply a common
test element by any nonzero element of $\mathfrak b^r$, absorbing this
bounded difference without changing the trace-generated test ideal.
This argument is local and applies with the same fixed boundary $\Delta$.

\begin{lemma}[Relative generation]\label{lem:gen}
There is an integer $c$, independent of $m$, such that
\[
 \mathscr F_m=t_m\otimes\OO_Y(mC-cF)
\]
is globally generated for every $m\geq1$.
\end{lemma}
\begin{proof}
Write $H=-F$, so $A=\OO(H)$. Choose $c_0$ with
$(c_0-i)H-L$ ample for every $1\leq i\leq\dim Y$.
Fix $m$ and an $\ell$ attaining $t_m$. Set $\lambda=m/\ell$.
The trace description of the big test ideal, in the version with
$\lceil\lambda p^e\rceil$, gives, for a sufficiently large effective
Cartier test divisor $D_0$ and any lower bound $e_0$,
\[
 t_m=\sum_{e\geq e_0}\operatorname{Image}\left[
 F_*^e\bigl(\mathfrak b_\ell^{\lceil\lambda p^e\rceil}
       \cdot\OO_Y(K_Y-p^eL-D_0)\bigr)\longrightarrow\OO_Y\right].
\]
The scheme version is obtained from the affine trace formula by a finite
cover and one Cartier divisor dominating its local test divisors. An ample
twist of the coherent test ideal supplies a nonzero global section for this
domination. The divisor $D_0$ may depend on $(m,\ell)$; its dependence
does not affect the choice of $c_0$.

Put $q=p^e$, $a_e=\lceil qm/\ell\rceil$, and
$j_e=\ell a_e-qm$. Then $0\leq j_e<\ell$.
The globally generated ideal
$\mathfrak b_\ell^{a_e}\OO(a_e\ell C)$ is generated by finite products
of global generators of $I_\ell$. After twisting the trace target by
$\OO(mC+c_0H)$ and using the projection formula, its source is a quotient
of finitely many copies of
\[
 Q_{e,j_e}=F_*^e\OO_Y\bigl(q(c_0H-L)+K_Y-D_0-j_eC\bigr).
\]
The quotient maps are multiplication by the indicated global section
products inside $F_*^e$. They are $\OO_Y$-linear for the Frobenius module
structure. This avoids pulling a base-module factor incorrectly through
Frobenius.

For $1\leq i\leq\dim Y$, the cohomology of $Q_{e,j}\otimes A^{-i}$ is
the cohomology, as an abelian group, of
\[
 \OO_Y(K_Y-D_0-jC)\otimes
 \OO_Y\bigl(q((c_0-i)H-L)\bigr).
\]
The first factor is a fixed coherent reflexive sheaf, even if $K_Y$ is
not Cartier. The second is a power of an ample Cartier line bundle.
Relative Serre vanishing over the affine base makes these groups zero
for all sufficiently large $e$, simultaneously for the finitely many
$j\in\{0,\ldots,\ell-1\}$ and $i$. Relative Castelnuovo--Mumford
regularity therefore makes $Q_{e,j}$ globally generated. The relative
regularity theorem used here applies to projective schemes over noetherian
affine bases, not only to varieties over a field.

All sufficiently large trace images, and hence their finite stabilized
sum, are globally generated. This proves generation with $c_0$ for the
fixed $m$. Only the Serre threshold depends on $m$; $c_0$ is fixed.
Finally replace $c_0$ by $c=c_0+1$ if needed. Tensoring by the globally
generated $A$ preserves generation and ensures that the section ideal
below is proper.
\end{proof}

Set $J_m=\Gamma(Y,\mathscr F_m)$. Hartogs again makes this an actual ideal
in $B$, and generation identifies $J_m\OO_Y=\mathscr F_m$.
Because $\Delta$, $C$, and $F$ are exceptional and $Y$ is regular off the
fiber, $t_m=\OO$ there. Thus $J_m$ is unit off $x$. The extra factor $A$
makes it proper: before tensoring, the generated sheaf has global sections
in $B$ and hence is a subsheaf of $\OO_Y$; afterwards it is contained in
$J\OO_Y$. Coherence and full-fiber support give a primary power contained
in $J_m$. Therefore
\[
 g_m=-\frac{v(J_m)}m
     =f_C-\frac{c\,v(J)+v(t_m)}m\in\CC,
 \qquad g_m\leq f_C.
\]

\section{Jacobian corrections vanish uniformly}
\begin{lemma}\label{lem:estimate}
On the entire normalized semivaluation section,
\[
 0\leq P f_C-g_m
 \leq\frac{c\,v(J)+v(\Jac(Y/k))+v(bF)}m
 \leq\frac{(c+a+b)N_J}m.
\]
\end{lemma}
\begin{proof}
Let $T_0=\tau(Y,\Delta)$. This nonzero coherent ideal is unit off the
fiber and contains $J^r\OO_Y$ for some $r$, so its valuation is finite
at every normalized semivaluation.
The elementary multiplication property of big test ideals gives
\[
 \mathfrak b_{m l}T_0\subset t_{m l}\subset s_{m l}.
\]
For completeness, use a common test element $z$ for $\tau(Y,\Delta)$
and $\tau(Y,\Delta,\mathfrak b_{ml})$ in an affine chart.
For $f\in\mathfrak b_{ml}$ and a Cartier map $\varphi$,
$f\varphi(F_*^e z)=\varphi(F_*^e(f^{p^e}z))$.
The latter is in the test ideal of exponent one by its trace formula.
Localization gives the displayed sheaf inclusion, and the term
$\ell=ml$ embeds it into the asymptotic ideal.

Jacobian-corrected subadditivity on reduced equidimensional algebras over
a perfect field, localized on $Y$, gives
\[
 \Jac(Y/k)^{l-1}s_{ml}\subset s_m^l.
\]
This is the big-test-ideal formula of Takagi and Smolkin. For asymptotic
ideals choose one common index attaining $s_m$ and $s_{ml}$ and iterate
the two-factor formula. There is no assertion of uncorrected singular
subadditivity.

Put $h=f_C-Pf_C$. By Lemma~\ref{lem:raw},
$h=\lim_l v(\mathfrak b_{ml})/(ml)$. The two inclusions imply
\[
 \frac{v(s_m)}m
 \leq\frac{v(s_{ml})}{ml}
       +\frac{l-1}{ml}v(\Jac(Y/k))
 \leq\frac{v(\mathfrak b_{ml})+v(T_0)}{ml}
       +\frac{l-1}{ml}v(\Jac(Y/k)).
\]
All these values are finite because the indicated ideals contain primary
powers. Let $l\to\infty$ at the fixed semivaluation to obtain
$v(s_m)\leq mh+v(\Jac(Y/k))$.

Since $0\leq\Delta\leq bF$, monotonicity in the boundary and the Cartier
shift formula give
\[
 \OO_Y(-bF)s_m\subset t_m\subset s_m,
 \qquad v(t_m)\leq v(s_m)+v(bF).
\]
Substitute into the formula for $g_m$. The lower bound follows from
$g_m\in\CC$ and $g_m\leq f_C$. The upper bound uses the fixed primary
power inclusions established after Lemma~\ref{lem:flatten}; no global
bound on arbitrary nonprimary ideals is assumed.
\end{proof}

Uniform convergence of $g_m$ now proves $Pf_C\in\CC$. This is recognition
by a uniform limit of actual geometric primary profiles, not by an increasing
supremum or a finite energy/capacity diagnostic.

\section{Continuous and upper semicontinuous obstacles}
The differences of elementary primary profiles form a rational vector
lattice containing constants. For instance,
$\max(g_1-h_1,g_2-h_2)
=\max(g_1+h_2,g_2+h_1)-(h_1+h_2)$, and sums and maxima of profiles are
represented, after common denominator clearing, by products and sums of
primary ideals. These functions separate normalized semivaluations:
the primary ideals $(f)+\mathfrak a^N$ record
$\min\{v(f),N\}$ for every $f$ and $N$. Their values distinguish finite
values and infinite values. The lattice Stone--Weierstrass theorem on the
compact section therefore gives uniform density of these signed model
functions in all finite continuous functions. A difference of two profiles
is represented by a supported Cartier divisor after principalizing their
primary ideals. Lemma~\ref{lem:flatten} then supplies the normal actual
primary model required above.

The envelope is nonexpansive without an unrestricted constant-translation
identity. If $\|\phi-\psi\|\leq\varepsilon$, any admissible $h$ below $\phi$
has $h-\varepsilon\in\DD$ below $\psi$. Thus
$P\phi\leq P\psi+\varepsilon$, and the reverse inequality is identical.
For continuous $\phi$, choose signed rational model functions $f_j\to\phi$
uniformly. Their envelopes belong to $\CC$ by Lemma~\ref{lem:estimate};
nonexpansiveness and closure give $P\phi\in\CC$. Raw maximality and the
geometric comparison give Theorem~\ref{thm:main} for continuous obstacles.
Positive homogeneity of the cone and the raw supremum gives
$P(tf)=tP(f)$ for every $t>0$. Thus the integral Cartier argument also
applies after denominator clearing to the rational signed model obstacles
used in this density step.

For a bounded upper semicontinuous $\phi$ on the compact Hausdorff section,
take the directed net of continuous majorants, with pointwise infimum
$\phi$. The envelopes of these majorants form one decreasing net in
$\CC$. Its limit belongs to $\DD$, is below $\phi$, and dominates every
$\DD$ minorant, by the continuous-obstacle result. This establishes the
claimed greatest minorant directly from one net. It does not rely on an
unproved closure under decreasing nets of arbitrary $\DD$ functions.

Primary-gauge changes are inherited on the same homogeneous ray space:
restriction to a new section divides both obstacle and admissible functions
by the same positive continuous degree-zero gauge ratio. Their order and
the raw supremum are preserved. This yields the exact envelope cocycle.
No normalization/conductor rule outside the isolated normal sector is added.

\section{Source scope and remaining original problem}
The applied external algebra is the current Stacks Fitting strict-transform
argument (15.27.3--4 and 15.8.4, 7--9), big test ideals and their trace,
localization, and Cartier-shift formulas in corrected CEMS v4 (2.2--4),
and the Jacobian subadditivity formula in corrected Smolkin v3 (3.11--12).
Takagi v2, 2.5--7 and 4.2--5, independently records the singular and
asymptotic formulas. GJKM 4.6 identifies the precise relative regularity
input over an affine base, referring to Brodmann--Sharp 20.4.13.
The new relative generation lemma is proved here; the projective-over-a-field
generation statements of Schwede or CEMS 4.5 are not substituted for it.
The CEMS finite-generation hypothesis on the anticanonical algebra is
not assumed for $Y$; the Cartier divisor $K_Y+\Delta$ is explicitly built.

There is a point to check carefully in the printed finite-field reduction
in Smolkin 3.12: a $k$-linear retraction from a perfect extension field is
not a ring homomorphism defining an extension-field module structure on
$k$. A valid descent of the lifting equation instead uses faithful flatness.
For a perfect extension $k\subset k'$, finite presentations and perfection
identify the base changes of the Frobenius Hom modules. Compatibility with
the diagonal is the kernel of a map between these modules, and restriction
to the quotient is another module map. Tensoring by $k'$ preserves these
kernels and images. Membership of $x\varphi$ in the image after extension
therefore descends. Passing to an infinite perfect extension proves the
finite-perfect-field case from the infinite case without treating a linear
retraction as multiplicative. This correction is part of the working proof
to be independently checked, not an attribution to the printed text.

The Hacon--Lamarche--Schwede published 2024 paper proves useful global
generation statements over complete local mixed-characteristic bases;
its 7.1 additionally assumes regularity, and 8.3 states global
subadditivity as a conjecture. Those statements do not supply a
mixed-characteristic extension of this argument. Nonisolated singularities
also leave the bounded correction step unresolved: their Jacobian and test
ideal may have nonprimary cosupport and unbounded values on the whole link.
Replacing a correction ideal by its primary truncation changes the direction
of the required containment and gives no solution.

The literal P12.3 requests geometric envelopes for excellent singular and
nonnormal germs, appropriate regularity, smooth specialization, gauge and
normalization/conductor rules, and the reduced-curve comparison. This
working paper proposes a new isolated normal algebraic positive-characteristic
sector. Its full proof must undergo independent verification. It is not an
original-problem closure and changes neither the solved count nor the
accepted-paper count. Primary source bodies, literal units, and source-page
pixels remain local; citations and original research prose may be checkpointed.

\begin{thebibliography}{9}
\bibitem{CEMS} A. Chiecchio, F. Enescu, L. E. Miller, K. Schwede,
\emph{Test ideals in rings with finitely generated anti-canonical algebras},
arXiv:1412.6453v4 (corrected 31 March 2016).
\bibitem{Smolkin} D. Smolkin, \emph{A new subadditivity formula for test ideals},
arXiv:1805.08739v3 (corrected 5 February 2019).
\bibitem{Takagi} S. Takagi, \emph{Formulas for multiplier ideals on singular
varieties}, arXiv:math/0410612v2 (31 January 2005).
\bibitem{GJKM} W. Gubler, P. Jell, K. K\"unnemann, F. Martin,
\emph{Continuity of plurisubharmonic envelopes in non-Archimedean geometry
and test ideals}, Section 4, especially Theorem 4.6.
\bibitem{BS} M. P. Brodmann, R. Y. Sharp,
\emph{Local Cohomology}, relative regularity theorem 20.4.13,
as identified by GJKM 4.6; full book proof not freshly inspected here.
\bibitem{HLS} C. Hacon, A. Lamarche, K. Schwede,
\emph{Global generation of test ideals in mixed characteristic and applications},
Algebraic Geometry 11 (5) (2024), 676--711.
\bibitem{Stacks} The Stacks Project, current tags
0535, 07ZA, 07ZC, 07ZD, 080Z, 07Z3, 080R, 07Z7.
\bibitem{G331} Qualified independently verified P12.2 geometric comparison
paper, continuation 331, with its mandatory source and reader-scope
qualification. Only its actual primary geometric comparison is used.
\end{thebibliography}
\end{document}
